%% notes-tex/database/spell-publications/sections/2-knockout.tex
%% [[experiments.015-spell.publications]]
%% https://github.com/Mjvolk3/torchcell/tree/main/notes-tex/database/spell-publications/sections/2-knockout.tex
%% SOURCE: numbers and tables from
%% experiments/015-spell/scripts/spell_publication_tables.py, reading
%% experiments/015-spell/results/spell_knockout_coverage.csv; prose hand-authored

\section{Gene-deletion coverage\secstatus{todo}}\label{sec:knockout}

A SPELL condition is described only by the free text of its column header, so
deletion coverage is estimated from that text. A header counts as a
\term{deletion condition} when a deletion marker is attached to a token that
resolves to a yeast gene. The markers are \texttt{-del}, a delta sign, the words
\texttt{delta}, \texttt{deletion}, \texttt{knockout}, \texttt{null} and
\texttt{KO}, and a lowercase gene name followed by a capital \texttt{D}. A token
resolves when it is a systematic name, a standard name, or an alias of exactly
one gene in the SGD R64-4-1 annotation.

By that rule \spellNKoStudies{} of the \spellNStudies{} studies have at least
one deletion condition, \spellNKoConditions{} of the \spellNConditions{}
conditions are deletion conditions, and \spellNKoGenes{} distinct genes are
deleted across the archive (Table~\ref{tab:knockout}). Coverage is concentrated:
\spellNKoStudiesTenPlus{} studies delete ten or more genes. Of the
\spellNKoGenes{} genes, \spellNKoGenesNew{} are absent from the 1,484 mutants of
Kemmeren 2014 Table~S1, the deletion compendium already built, and they are
spread over \spellNKoStudiesNew{} studies.

The estimate errs in both directions. It misses a deletion written without a
marker, such as \texttt{ptc1 vs wt}, and it counts a marked allele that is not a
full deletion, such as the stabilized \texttt{SIC1$\Delta$3P}. Every call is
listed with its header in
\file{experiments/015-spell/results/spell_knockout_conditions.csv}. No call was
checked against its paper.

\subsection{Genotypes new to the built datasets\secstatus{todo}}\label{sec:new}

Kemmeren 2014 (1,484 single deletions) and Sameith 2015 (single and double
deletions) are the built expression datasets with a gene perturbation. A
single-gene genotype counts as covered when its gene is in Kemmeren 2014
Table~S1, and any genotype counts as covered when it occurs in the Sameith 2015
folder of the archive, which stands in for the built Sameith datasets. A
genotype that is neither is \term{new}.

\spellNNewStudies{} studies carry at least one new genotype, for
\spellNNewGenotypes{} genotypes in all. \spellNNewKoGenotypes{} of those are
deletion genotypes spread over \spellNNewKoStudies{} studies, and
\spellNNewNamedGenotypes{} are gene-named genotypes
(Sec.~\ref{sec:gene-named}), nearly all in Hughes 2000 and Mnaimneh 2004. Only
\spellNNewStudiesFivePlus{} studies add five or more. The same genotype can be
counted in two studies, so the per-study numbers sum to more than the number of
distinct new genotypes.

%% GENERATED FILE -- do not hand-edit.
%% SOURCE: experiments/015-spell/scripts/spell_publication_tables.py
\begin{landscape}
\begingroup
\footnotesize
\setlength{\tabcolsep}{4pt}
\renewcommand{\arraystretch}{1.1}
\begin{longtable}{@{}r@{\hspace{5pt}} L{128mm} r r r r r L{20mm} l@{}}
\caption[]{Studies with at least one deletion call, most deleted genes first. \emph{\#} is the row number in Table~\ref{tab:all}. \emph{KO genes} is the number of distinct genes a deletion marker is attached to in the study's column headers. \emph{genotypes} counts distinct sets of deleted genes, so a double mutant is one genotype, and \emph{new} is how many of those genotypes are in neither Kemmeren 2014 nor Sameith 2015 (Sec.~\ref{sec:new}); a dash is zero. \emph{KO cond.} is the number of conditions with a call and \emph{cond.} the study's total. The calls are read from header text and are estimates (Sec.~\ref{sec:knockout}).}
\label{tab:knockout}\\
\toprule
\textbf{\#} & \textbf{Publication} & \textbf{KO genes} & \textbf{new} & \textbf{genotypes} & \textbf{KO cond.} & \textbf{cond.} & \textbf{GEO} & \textbf{link} \\
\midrule
\endfirsthead
\multicolumn{9}{@{}l}{\footnotesize\emph{Table~\ref{tab:knockout}, continued}}\\
\toprule
\textbf{\#} & \textbf{Publication} & \textbf{KO genes} & \textbf{new} & \textbf{genotypes} & \textbf{KO cond.} & \textbf{cond.} & \textbf{GEO} & \textbf{link} \\
\midrule
\endhead
\bottomrule
\endfoot
4 & \textbf{Hu 2007}\newline \sourcetext{Genetic reconstruction of a functional transcriptional regulatory network. \emph{Nat Genet}} 2007. & 269 & 14 & 269 & 269 & 269 & \href{https://www.ncbi.nlm.nih.gov/geo/query/acc.cgi?acc=GSE4654}{GSE4654} & \href{https://doi.org/10.1038/ng2012}{doi} \\
\addlinespace[3pt]
51 & \textbf{Lenstra 2011}\newline \sourcetext{The specificity and topology of chromatin interaction pathways in yeast. \emph{Mol Cell}} 2011. & 162 & -- & 162 & 313 & 313 & \href{https://www.ncbi.nlm.nih.gov/geo/query/acc.cgi?acc=GSE25909}{GSE25909} & \href{https://doi.org/10.1016/j.molcel.2011.03.026}{doi} \\
\addlinespace[3pt]
3 & \textbf{van Wageningen 2010}\newline \sourcetext{Functional overlap and regulatory links shape genetic interactions between signaling pathways. \emph{Cell}} 2010. & 133 & 14 & 145 & 290 & 377 & \href{https://www.ncbi.nlm.nih.gov/geo/query/acc.cgi?acc=GSE25644}{GSE25644} & \href{https://doi.org/10.1016/j.cell.2010.11.021}{doi} \\
\addlinespace[3pt]
8 & \textbf{Apweiler 2012}\newline \sourcetext{Yeast glucose pathways converge on the transcriptional regulation of trehalose biosynthesis. \emph{BMC Genomics}} 2012. & 91 & 3 & 93 & 183 & 275 & \href{https://www.ncbi.nlm.nih.gov/geo/query/acc.cgi?acc=GSE33098}{GSE33098}\newline \href{https://www.ncbi.nlm.nih.gov/geo/query/acc.cgi?acc=GSE33097}{GSE33097} & \href{https://doi.org/10.1186/1471-2164-13-239}{doi} \\
\addlinespace[3pt]
52 & \textbf{Sameith 2015}\newline \sourcetext{A high-resolution gene expression atlas of epistasis between gene-specific transcription factors exposes potential mechanisms for genetic interactions. \emph{BMC Biol}} 2015. & 82 & -- & 154 & 287 & 287 & \href{https://www.ncbi.nlm.nih.gov/geo/query/acc.cgi?acc=GSE42536}{GSE42536} & \href{https://doi.org/10.1186/s12915-015-0222-5}{doi} \\
\addlinespace[3pt]
53 & \textbf{Chua 2006}\newline \sourcetext{Identifying transcription factor functions and targets by phenotypic activation. \emph{Proc Natl Acad Sci U S A}} 2006. & 51 & -- & 51 & 102 & 270 & \href{https://www.ncbi.nlm.nih.gov/geo/query/acc.cgi?acc=GSE5499}{GSE5499} & \href{https://doi.org/10.1073/pnas.0605140103}{doi} \\
\addlinespace[3pt]
128 & \textbf{van de Pasch 2013}\newline \sourcetext{Centromere binding and a conserved role in chromosome stability for SUMO-dependent ubiquitin ligases. \emph{PLoS One}} 2013. & 17 & -- & 17 & 34 & 34 & \href{https://www.ncbi.nlm.nih.gov/geo/query/acc.cgi?acc=GSE33929}{GSE33929} & \href{https://doi.org/10.1371/journal.pone.0065628}{doi} \\
\addlinespace[3pt]
132 & \textbf{Garcia-Oliver 2013}\newline \sourcetext{A novel role for Sem1 and TREX-2 in transcription involves their impact on recruitment and H2B deubiquitylation activity of SAGA. \emph{Nucleic Acids Res}} 2013. & 16 & -- & 16 & 32 & 32 & \href{https://www.ncbi.nlm.nih.gov/geo/query/acc.cgi?acc=GSE39861}{GSE39861} & \href{https://doi.org/10.1093/nar/gkt272}{doi} \\
\addlinespace[3pt]
29 & \textbf{Huang 2002}\newline \sourcetext{Dissection of a complex phenotype by functional genomics reveals roles for the yeast cyclin-dependent protein kinase Pho85 in stress adaptation and cell integrity. \emph{Mol Cell Biol}} 2002. & 12 & 1 & 12 & 12 & 20 & \textbf{none} & \href{https://doi.org/10.1128/MCB.22.14.5076-5088.2002}{doi} \\
\addlinespace[3pt]
188 & \textbf{Lenstra 2013}\newline \sourcetext{The role of Ctk1 kinase in termination of small non-coding RNAs. \emph{PLoS One}} 2013. & 7 & -- & 7 & 16 & 22 & \href{https://www.ncbi.nlm.nih.gov/geo/query/acc.cgi?acc=GSE40254}{GSE40254} & \href{https://doi.org/10.1371/journal.pone.0080495}{doi} \\
\addlinespace[3pt]
25 & \textbf{Munding 2010}\newline \sourcetext{Integration of a splicing regulatory network within the meiotic gene expression program of Saccharomyces cerevisiae. \emph{Genes Dev}} 2010. & 6 & 1 & 6 & 20 & 48 & \href{https://www.ncbi.nlm.nih.gov/geo/query/acc.cgi?acc=GSE24675}{GSE24675}\newline \href{https://www.ncbi.nlm.nih.gov/geo/query/acc.cgi?acc=GSE24681}{GSE24681} & \href{https://doi.org/10.1101/gad.1977410}{doi} \\
\addlinespace[3pt]
15 & \textbf{Gitter 2013}\newline \sourcetext{Linking the signaling cascades and dynamic regulatory networks controlling stress responses. \emph{Genome Res}} 2013. & 6 & 2 & 6 & 12 & 16 & \href{https://www.ncbi.nlm.nih.gov/geo/query/acc.cgi?acc=GSE28213}{GSE28213} & \href{https://doi.org/10.1101/gr.138628.112}{doi} \\
\addlinespace[3pt]
5 & \textbf{Zhou 2011}\newline \sourcetext{Integrated approaches reveal determinants of genome-wide binding and function of the transcription factor Pho4. \emph{Mol Cell}} 2011. & 5 & 4 & 5 & 14 & 66 & \href{https://www.ncbi.nlm.nih.gov/geo/query/acc.cgi?acc=GSE23580}{GSE23580} & \href{https://doi.org/10.1016/j.molcel.2011.05.025}{doi} \\
\addlinespace[3pt]
355 & \textbf{Wan 2011}\newline \sourcetext{Transcriptome profiling reveals a novel role for trichostatin A in antagonizing histone chaperone Chz1 mediated telomere anti-silencing. \emph{FEBS Lett}} 2011. & 5 & -- & 5 & 10 & 11 & \href{https://www.ncbi.nlm.nih.gov/geo/query/acc.cgi?acc=GSE30228}{GSE30228} & \href{https://doi.org/10.1016/j.febslet.2011.06.036}{doi} \\
\addlinespace[3pt]
13 & \textbf{Cho 2019}\newline \sourcetext{The cell-cycle transcriptional network generates and transmits a pulse of transcription once each cell cycle. \emph{Cell Cycle}} 2019. & 4 & 2 & 2 & 96 & 270 & \href{https://www.ncbi.nlm.nih.gov/geo/query/acc.cgi?acc=GSE75694}{GSE75694} & \href{https://doi.org/10.1080/15384101.2019.1570655}{doi} \\
\addlinespace[3pt]
6 & \textbf{Braberg 2013}\newline \sourcetext{From structure to systems: high-resolution, quantitative genetic analysis of RNA polymerase II. \emph{Cell}} 2013. & 4 & 4 & 4 & 55 & 57 & \href{https://www.ncbi.nlm.nih.gov/geo/query/acc.cgi?acc=GSE47429}{GSE47429} & \href{https://doi.org/10.1016/j.cell.2013.07.033}{doi} \\
\addlinespace[3pt]
9 & \textbf{Morillo-Huesca 2010}\newline \sourcetext{The SWR1 histone replacement complex causes genetic instability and genome-wide transcription misregulation in the absence of H2A.Z. \emph{PLoS One}} 2010. & 4 & 3 & 7 & 21 & 24 & \href{https://www.ncbi.nlm.nih.gov/geo/query/acc.cgi?acc=GSE21571}{GSE21571} & \href{https://doi.org/10.1371/journal.pone.0012143}{doi} \\
\addlinespace[3pt]
16 & \textbf{Li 2010}\newline \sourcetext{Thiamine biosynthesis in Saccharomyces cerevisiae is regulated by the NAD+-dependent histone deacetylase Hst1. \emph{Mol Cell Biol}} 2010. & 4 & 2 & 4 & 8 & 14 & \href{https://www.ncbi.nlm.nih.gov/geo/query/acc.cgi?acc=GSE18488}{GSE18488} & \href{https://doi.org/10.1128/MCB.01590-09}{doi} \\
\addlinespace[3pt]
525 & \textbf{Baetz 2001}\newline \sourcetext{Transcriptional coregulation by the cell integrity mitogen-activated protein kinase Slt2 and the cell cycle regulator Swi4. \emph{Mol Cell Biol}} 2001. & 4 & -- & 4 & 5 & 5 & \textbf{none} & \href{https://doi.org/10.1128/MCB.21.19.6515-6528.2001}{doi} \\
\addlinespace[3pt]
368 & \textbf{Moxley 2009}\newline \sourcetext{Linking high-resolution metabolic flux phenotypes and transcriptional regulation in yeast modulated by the global regulator Gcn4p. \emph{Proc Natl Acad Sci U S A}} 2009. & 4 & -- & 4 & 4 & 10 & \href{https://www.ncbi.nlm.nih.gov/geo/query/acc.cgi?acc=GSE4709}{GSE4709}\newline \href{https://www.ncbi.nlm.nih.gov/geo/query/acc.cgi?acc=GSE7369}{GSE7369} & \href{https://doi.org/10.1073/pnas.0811091106}{doi} \\
\addlinespace[3pt]
76 & \textbf{Capaldi 2008}\newline \sourcetext{Structure and function of a transcriptional network activated by the MAPK Hog1. \emph{Nat Genet}} 2008. & 3 & -- & 3 & 38 & 85 & \href{https://www.ncbi.nlm.nih.gov/geo/query/acc.cgi?acc=GSE12270}{GSE12270} & \href{https://doi.org/10.1038/ng.235}{doi} \\
\addlinespace[3pt]
7 & \textbf{Verzijlbergen 2011}\newline \sourcetext{A barcode screen for epigenetic regulators reveals a role for the NuB4/HAT-B histone acetyltransferase complex in histone turnover. \emph{PLoS Genet}} 2011. & 3 & 4 & 7 & 24 & 36 & \href{https://www.ncbi.nlm.nih.gov/geo/query/acc.cgi?acc=GSE30168}{GSE30168} & \href{https://doi.org/10.1371/journal.pgen.1002284}{doi} \\
\addlinespace[3pt]
19 & \textbf{Mutiu 2007}\newline \sourcetext{The role of histone ubiquitylation and deubiquitylation in gene expression as determined by the analysis of an HTB1(K123R) Saccharomyces cerevisiae strain. \emph{Mol Genet Genomics}} 2007. & 3 & 2 & 4 & 9 & 9 & \href{https://www.ncbi.nlm.nih.gov/geo/query/acc.cgi?acc=GSE6316}{GSE6316} & \href{https://doi.org/10.1007/s00438-007-0212-6}{doi} \\
\addlinespace[3pt]
260 & \textbf{Santos-Pereira 2014}\newline \sourcetext{A genome-wide function of THSC/TREX-2 at active genes prevents transcription-replication collisions. \emph{Nucleic Acids Res}} 2014. & 3 & -- & 3 & 9 & 15 & \href{https://www.ncbi.nlm.nih.gov/geo/query/acc.cgi?acc=GSE56702}{GSE56702} & \href{https://doi.org/10.1093/nar/gku906}{doi} \\
\addlinespace[3pt]
10 & \textbf{Thibault 2011}\newline \sourcetext{The unfolded protein response supports cellular robustness as a broad-spectrum compensatory pathway. \emph{Proc Natl Acad Sci U S A}} 2011. & 3 & 3 & 3 & 9 & 15 & \href{https://www.ncbi.nlm.nih.gov/geo/query/acc.cgi?acc=GSE33844}{GSE33844} & \href{https://doi.org/10.1073/pnas.1117184109}{doi} \\
\addlinespace[3pt]
30 & \textbf{Thibault 2012}\newline \sourcetext{The membrane stress response buffers lethal effects of lipid disequilibrium by reprogramming the protein homeostasis network. \emph{Mol Cell}} 2012. & 3 & 1 & 3 & 9 & 16 & \href{https://www.ncbi.nlm.nih.gov/geo/query/acc.cgi?acc=GSE39419}{GSE39419} & \href{https://doi.org/10.1016/j.molcel.2012.08.016}{doi} \\
\addlinespace[3pt]
11 & \textbf{Gomez-Gonzalez 2011}\newline \sourcetext{Genome-wide function of THO/TREX in active genes prevents R-loop-dependent replication obstacles. \emph{EMBO J}} 2011. & 3 & 3 & 3 & 7 & 10 & \href{https://www.ncbi.nlm.nih.gov/geo/query/acc.cgi?acc=GSE24802}{GSE24802} & \href{https://doi.org/10.1038/emboj.2011.206}{doi} \\
\addlinespace[3pt]
197 & \textbf{Smith 2007}\newline \sourcetext{Transcriptional responses to fatty acid are coordinated by combinatorial control. \emph{Mol Syst Biol}} 2007. & 3 & -- & 3 & 6 & 20 & \href{https://www.ncbi.nlm.nih.gov/geo/query/acc.cgi?acc=GSE5862}{GSE5862} & \href{https://doi.org/10.1038/msb4100157}{doi} \\
\addlinespace[3pt]
463 & \textbf{Van de Vosse 2013}\newline \sourcetext{A role for the nucleoporin Nup170p in chromatin structure and gene silencing. \emph{Cell}} 2013. & 3 & -- & 3 & 3 & 6 & \href{https://www.ncbi.nlm.nih.gov/geo/query/acc.cgi?acc=GSE36793}{GSE36793} & \href{https://doi.org/10.1016/j.cell.2013.01.049}{doi} \\
\addlinespace[3pt]
110 & \textbf{Gardner 2005}\newline \sourcetext{Ubp10/Dot4p regulates the persistence of ubiquitinated histone H2B: distinct roles in telomeric silencing and general chromatin. \emph{Mol Cell Biol}} 2005. & 2 & -- & 2 & 18 & 42 & \href{https://www.ncbi.nlm.nih.gov/geo/query/acc.cgi?acc=GSE2330}{GSE2330}\newline \href{https://www.ncbi.nlm.nih.gov/geo/query/acc.cgi?acc=GSE2329}{GSE2329} & \href{https://doi.org/10.1128/MCB.25.14.6123-6139.2005}{doi} \\
\addlinespace[3pt]
163 & \textbf{Zhu 2009}\newline \sourcetext{High-resolution DNA-binding specificity analysis of yeast transcription factors. \emph{Genome Res}} 2009. & 2 & -- & 3 & 18 & 24 & \href{https://www.ncbi.nlm.nih.gov/geo/query/acc.cgi?acc=GSE13684}{GSE13684} & \href{https://doi.org/10.1101/gr.090233.108}{doi} \\
\addlinespace[3pt]
27 & \textbf{Ansari 2012}\newline \sourcetext{Distinct role of Mediator tail module in regulation of SAGA-dependent, TATA-containing genes in yeast. \emph{EMBO J}} 2012. & 2 & 1 & 2 & 12 & 21 & \href{https://www.ncbi.nlm.nih.gov/geo/query/acc.cgi?acc=GSE31774}{GSE31774} & \href{https://doi.org/10.1038/emboj.2011.362}{doi} \\
\addlinespace[3pt]
406 & \textbf{Dowell 2010}\newline \sourcetext{Chromatin-dependent binding of the S. cerevisiae HMGB protein Nhp6A affects nucleosome dynamics and transcription. \emph{Genes Dev}} 2010. & 2 & -- & 2 & 8 & 8 & \href{https://www.ncbi.nlm.nih.gov/geo/query/acc.cgi?acc=GSE23607}{GSE23607} & \href{https://doi.org/10.1101/gad.1948910}{doi} \\
\addlinespace[3pt]
28 & \textbf{Borde 2009}\newline \sourcetext{Histone H3 lysine 4 trimethylation marks meiotic recombination initiation sites. \emph{EMBO J}} 2009. & 2 & 1 & 1 & 7 & 21 & \href{https://www.ncbi.nlm.nih.gov/geo/query/acc.cgi?acc=GSE10948}{GSE10948}\newline \href{https://www.ncbi.nlm.nih.gov/geo/query/acc.cgi?acc=GSE10944}{GSE10944}\newline \href{https://www.ncbi.nlm.nih.gov/geo/query/acc.cgi?acc=GSE10947}{GSE10947} & \href{https://doi.org/10.1038/emboj.2008.257}{doi} \\
\addlinespace[3pt]
20 & \textbf{Belli 2004}\newline \sourcetext{Saccharomyces cerevisiae glutaredoxin 5-deficient cells subjected to continuous oxidizing conditions are affected in the expression of specific sets of genes. \emph{J Biol Chem}} 2004. & 2 & 2 & 2 & 6 & 9 & \textbf{none} & \href{https://doi.org/10.1074/jbc.M311879200}{doi} \\
\addlinespace[3pt]
12 & \textbf{Cain 2012}\newline \sourcetext{A conserved transcriptional regulator governs fungal morphology in widely diverged species. \emph{Genetics}} 2012. & 2 & 3 & 3 & 6 & 8 & \href{https://www.ncbi.nlm.nih.gov/geo/query/acc.cgi?acc=GSE32550}{GSE32550} & \href{https://doi.org/10.1534/genetics.111.134080}{doi} \\
\addlinespace[3pt]
297 & \textbf{Carroll 2001}\newline \sourcetext{Chemical inhibition of the Pho85 cyclin-dependent kinase reveals a role in the environmental stress response. \emph{Proc Natl Acad Sci U S A}} 2001. & 2 & -- & 2 & 6 & 12 & \textbf{none} & \href{https://doi.org/10.1073/pnas.211195798}{doi} \\
\addlinespace[3pt]
191 & \textbf{Kunkel 2019}\newline \sourcetext{Integrated TORC1 and PKA signaling control the temporal activation of glucose-induced gene expression in yeast. \emph{Nat Commun}} 2019. & 2 & -- & 1 & 6 & 21 & \href{https://www.ncbi.nlm.nih.gov/geo/query/acc.cgi?acc=GSE133591}{GSE133591} & \href{https://doi.org/10.1038/s41467-019-11540-y}{doi} \\
\addlinespace[3pt]
21 & \textbf{Chen 2017}\newline \sourcetext{HAL2 overexpression induces iron acquisition in bdf1$\Delta$ cells and enhances their salt resistance. \emph{Curr Genet}} 2017. & 2 & 2 & 2 & 5 & 5 & \href{https://www.ncbi.nlm.nih.gov/geo/query/acc.cgi?acc=GSE75828}{GSE75828} & \href{https://doi.org/10.1007/s00294-016-0628-9}{doi} \\
\addlinespace[3pt]
38 & \textbf{Leadsham 2010}\newline \sourcetext{cAMP/PKA signaling balances respiratory activity with mitochondria dependent apoptosis via transcriptional regulation. \emph{BMC Cell Biol}} 2010. & 2 & 1 & 2 & 5 & 8 & \href{https://www.ncbi.nlm.nih.gov/geo/query/acc.cgi?acc=GSE25541}{GSE25541} & \href{https://doi.org/10.1186/1471-2121-11-92}{doi} \\
\addlinespace[3pt]
450 & \textbf{Chen 2013}\newline \sourcetext{Stabilization of the promoter nucleosomes in nucleosome-free regions by the yeast Cyc8-Tup1 corepressor. \emph{Genome Res}} 2013. & 2 & -- & 2 & 4 & 7 & \href{https://www.ncbi.nlm.nih.gov/geo/query/acc.cgi?acc=GSE37466}{GSE37466} & \href{https://doi.org/10.1101/gr.141952.112}{doi} \\
\addlinespace[3pt]
407 & \textbf{Kugou 2007}\newline \sourcetext{Mre11 mediates gene regulation in yeast spore development. \emph{Genes Genet Syst}} 2007. & 2 & -- & 2 & 4 & 8 & \href{https://www.ncbi.nlm.nih.gov/geo/query/acc.cgi?acc=GSE6620}{GSE6620} & \href{https://doi.org/10.1266/ggs.82.21}{doi} \\
\addlinespace[3pt]
46 & \textbf{Lai 2013}\newline \sourcetext{Yeast hEST1A/B (SMG5/6)-like proteins contribute to environment-sensing adaptive gene expression responses. \emph{G3 (Bethesda)}} 2013. & 2 & 1 & 1 & 4 & 4 & \href{https://www.ncbi.nlm.nih.gov/geo/query/acc.cgi?acc=GSE48956}{GSE48956} & \href{https://doi.org/10.1534/g3.113.006924}{doi} \\
\addlinespace[3pt]
298 & \textbf{Navlakha 2012}\newline \sourcetext{A network-based approach for predicting missing pathway interactions. \emph{PLoS Comput Biol}} 2012. & 2 & -- & 2 & 4 & 12 & \href{https://www.ncbi.nlm.nih.gov/geo/query/acc.cgi?acc=GSE28213}{GSE28213} & \href{https://doi.org/10.1371/journal.pcbi.1002640}{doi} \\
\addlinespace[3pt]
39 & \textbf{Revilleza 2022}\newline \sourcetext{Regulation of CLB6 expression by the cytoplasmic deadenylase Ccr4 through its coding and 3' UTR regions. \emph{PLoS One}} 2022. & 2 & 1 & 2 & 4 & 6 & \href{https://www.ncbi.nlm.nih.gov/geo/query/acc.cgi?acc=GSE198743}{GSE198743} & \href{https://doi.org/10.1371/journal.pone.0268283}{doi} \\
\addlinespace[3pt]
464 & \textbf{Tamaki 2005}\newline \sourcetext{Glucose-dependent cell size is regulated by a G protein-coupled receptor system in yeast Saccharomyces cerevisiae. \emph{Genes Cells}} 2005. & 2 & -- & 2 & 4 & 6 & \href{https://www.ncbi.nlm.nih.gov/geo/query/acc.cgi?acc=GSE970}{GSE970} & \href{https://doi.org/10.1111/j.1365-2443.2005.00828.x}{doi} \\
\addlinespace[3pt]
18 & \textbf{Hazelwood 2008}\newline \sourcetext{The Ehrlich pathway for fusel alcohol production: a century of research on Saccharomyces cerevisiae metabolism. \emph{Appl Environ Microbiol}} 2008. & 2 & 2 & 2 & 3 & 12 & \href{https://www.ncbi.nlm.nih.gov/geo/query/acc.cgi?acc=GSE9590}{GSE9590} & \href{https://doi.org/10.1128/AEM.02625-07}{doi} \\
\addlinespace[3pt]
356 & \textbf{Penheiter 2005}\newline \sourcetext{A posttranscriptional role for the yeast Paf1-RNA polymerase II complex is revealed by identification of primary targets. \emph{Mol Cell}} 2005. & 2 & -- & 2 & 3 & 11 & \href{https://www.ncbi.nlm.nih.gov/geo/query/acc.cgi?acc=GSE3200}{GSE3200} & \href{https://doi.org/10.1016/j.molcel.2005.08.023}{doi} \\
\addlinespace[3pt]
573 & \textbf{Byeon 2013}\newline \sourcetext{The ATP-dependent chromatin remodeling enzyme Fun30 represses transcription by sliding promoter-proximal nucleosomes. \emph{J Biol Chem}} 2013. & 2 & -- & 2 & 2 & 3 & \href{https://www.ncbi.nlm.nih.gov/geo/query/acc.cgi?acc=GSE48570}{GSE48570} & \href{https://doi.org/10.1074/jbc.M113.471979}{doi} \\
\addlinespace[3pt]
22 & \textbf{Johnson 2014}\newline \sourcetext{Methionine restriction activates the retrograde response and confers both stress tolerance and lifespan extension to yeast, mouse and human cells. \emph{PLoS One}} 2014. & 2 & 2 & 2 & 2 & 3 & \href{https://www.ncbi.nlm.nih.gov/geo/query/acc.cgi?acc=GSE56641}{GSE56641} & \href{https://doi.org/10.1371/journal.pone.0097729}{doi} \\
\addlinespace[3pt]
574 & \textbf{Kamata 2023}\newline \sourcetext{Spt3 and Spt8 Are Involved in the Formation of a Silencing Boundary by Interacting with TATA-Binding Protein. \emph{Biomolecules}} 2023. & 2 & -- & 2 & 2 & 3 & \href{https://www.ncbi.nlm.nih.gov/geo/query/acc.cgi?acc=GSE220290}{GSE220290} & \href{https://doi.org/10.3390/biom13040619}{doi} \\
\addlinespace[3pt]
534 & \textbf{Orlandi 2004}\newline \sourcetext{Transcriptional profiling of ubp10 null mutant reveals altered subtelomeric gene expression and insurgence of oxidative stress response. \emph{J Biol Chem}} 2004. & 2 & -- & 2 & 2 & 4 & \textbf{none} & \href{https://doi.org/10.1074/jbc.M306464200}{doi} \\
\addlinespace[3pt]
575 & \textbf{Uluisik 2011}\newline \sourcetext{Boron stress activates the general amino acid control mechanism and inhibits protein synthesis. \emph{PLoS One}} 2011. & 2 & -- & 2 & 2 & 3 & \href{https://www.ncbi.nlm.nih.gov/geo/query/acc.cgi?acc=GSE29887}{GSE29887} & \href{https://doi.org/10.1371/journal.pone.0027772}{doi} \\
\addlinespace[3pt]
23 & \textbf{Huisinga 2007}\newline \sourcetext{A TATA binding protein regulatory network that governs transcription complex assembly. \emph{Genome Biol}} 2007. & 1 & 1 & 1 & 63 & 149 & \href{https://www.ncbi.nlm.nih.gov/geo/query/acc.cgi?acc=GSE7385}{GSE7385} & \href{https://doi.org/10.1186/gb-2007-8-4-r46}{doi} \\
\addlinespace[3pt]
111 & \textbf{Johansson 2007}\newline \sourcetext{Association of yeast Upf1p with direct substrates of the NMD pathway. \emph{Proc Natl Acad Sci U S A}} 2007. & 1 & -- & 1 & 20 & 40 & \href{https://www.ncbi.nlm.nih.gov/geo/query/acc.cgi?acc=GSE9482}{GSE9482} & \href{https://doi.org/10.1073/pnas.0709257105}{doi} \\
\addlinespace[3pt]
31 & \textbf{Pitkanen 2004}\newline \sourcetext{Excess mannose limits the growth of phosphomannose isomerase PMI40 deletion strain of Saccharomyces cerevisiae. \emph{J Biol Chem}} 2004. & 1 & 1 & 1 & 15 & 15 & \textbf{none} & \href{https://doi.org/10.1074/jbc.M410619200}{doi} \\
\addlinespace[3pt]
154 & \textbf{Cipollina 2008}\newline \sourcetext{Saccharomyces cerevisiae SFP1: at the crossroads of central metabolism and ribosome biogenesis. \emph{Microbiology (Reading)}} 2008. & 1 & -- & 1 & 13 & 26 & \href{https://www.ncbi.nlm.nih.gov/geo/query/acc.cgi?acc=GSE9644}{GSE9644} & \href{https://doi.org/10.1099/mic.0.2008/017392-0}{doi} \\
\addlinespace[3pt]
32 & \textbf{Barbara 2007}\newline \sourcetext{The transcription factor Gcr1 stimulates cell growth by participating in nutrient-responsive gene expression on a global level. \emph{Mol Genet Genomics}} 2007. & 1 & 1 & 1 & 12 & 12 & \href{https://www.ncbi.nlm.nih.gov/geo/query/acc.cgi?acc=GSE5027}{GSE5027} & \href{https://doi.org/10.1007/s00438-006-0182-0}{doi} \\
\addlinespace[3pt]
33 & \textbf{Crosas 2015}\newline \sourcetext{The yeast $\zeta$-crystallin/NADPH:quinone oxidoreductase (Zta1p) is under nutritional control by the target of rapamycin pathway and is involved in the regulation of argininosuccinate lyase mRNA half-life. \emph{FEBS J}} 2015. & 1 & 1 & 1 & 12 & 12 & \href{https://www.ncbi.nlm.nih.gov/geo/query/acc.cgi?acc=GSE45297}{GSE45297} & \href{https://doi.org/10.1111/febs.13246}{doi} \\
\addlinespace[3pt]
138 & \textbf{Zhou 2015}\newline \sourcetext{UV induced ubiquitination of the yeast Rad4-Rad23 complex promotes survival by regulating cellular dNTP pools. \emph{Nucleic Acids Res}} 2015. & 1 & -- & 1 & 12 & 30 & \href{https://www.ncbi.nlm.nih.gov/geo/query/acc.cgi?acc=GSE23204}{GSE23204} & \href{https://doi.org/10.1093/nar/gkv680}{doi} \\
\addlinespace[3pt]
24 & \textbf{Levy 2007}\newline \sourcetext{Strategy of transcription regulation in the budding yeast. \emph{PLoS One}} 2007. & 1 & 1 & 1 & 10 & 95 & \href{https://www.ncbi.nlm.nih.gov/geo/query/acc.cgi?acc=GSE6302}{GSE6302} & \href{https://doi.org/10.1371/journal.pone.0000250}{doi} \\
\addlinespace[3pt]
299 & \textbf{Boer 2005}\newline \sourcetext{Contribution of the Saccharomyces cerevisiae transcriptional regulator Leu3p to physiology and gene expression in nitrogen- and carbon-limited chemostat cultures. \emph{FEMS Yeast Res}} 2005. & 1 & -- & 1 & 6 & 12 & \textbf{none} & \href{https://doi.org/10.1016/j.femsyr.2005.04.003}{doi} \\
\addlinespace[3pt]
208 & \textbf{Busti 2012}\newline \sourcetext{Overexpression of Far1, a cyclin-dependent kinase inhibitor, induces a large transcriptional reprogramming in which RNA synthesis senses Far1 in a Sfp1-mediated way. \emph{Biotechnol Adv}} 2012. & 1 & -- & 1 & 6 & 18 & \href{https://www.ncbi.nlm.nih.gov/geo/query/acc.cgi?acc=GSE31143}{GSE31143} & \href{https://doi.org/10.1016/j.biotechadv.2011.09.007}{doi} \\
\addlinespace[3pt]
301 & \textbf{Cipollina 2008}\newline \sourcetext{Revisiting the role of yeast Sfp1 in ribosome biogenesis and cell size control: a chemostat study. \emph{Microbiology (Reading)}} 2008. & 1 & -- & 1 & 6 & 12 & \href{https://www.ncbi.nlm.nih.gov/geo/query/acc.cgi?acc=GSE5238}{GSE5238} & \href{https://doi.org/10.1099/mic.0.2007/011767-0}{doi} \\
\addlinespace[3pt]
41 & \textbf{Gaillard 2019}\newline \sourcetext{The Nup84 complex coordinates the DNA damage response to warrant genome integrity. \emph{Nucleic Acids Res}} 2019. & 1 & 1 & 1 & 6 & 6 & \href{https://www.ncbi.nlm.nih.gov/geo/query/acc.cgi?acc=GSE113118}{GSE113118} & \href{https://doi.org/10.1093/nar/gkz066}{doi} \\
\addlinespace[3pt]
357 & \textbf{Guan 2012}\newline \sourcetext{Cellular memory of acquired stress resistance in Saccharomyces cerevisiae. \emph{Genetics}} 2012. & 1 & -- & 1 & 6 & 11 & \href{https://www.ncbi.nlm.nih.gov/geo/query/acc.cgi?acc=GSE32196}{GSE32196} & \href{https://doi.org/10.1534/genetics.112.143016}{doi} \\
\addlinespace[3pt]
303 & \textbf{Hou 2014}\newline \sourcetext{Management of the endoplasmic reticulum stress by activation of the heat shock response in yeast. \emph{FEMS Yeast Res}} 2014. & 1 & -- & 1 & 6 & 12 & \href{https://www.ncbi.nlm.nih.gov/geo/query/acc.cgi?acc=GSE39311}{GSE39311} & \href{https://doi.org/10.1111/1567-1364.12125}{doi} \\
\addlinespace[3pt]
304 & \textbf{Pedersen 2012}\newline \sourcetext{DNA Topoisomerases maintain promoters in a state competent for transcriptional activation in Saccharomyces cerevisiae. \emph{PLoS Genet}} 2012. & 1 & -- & 1 & 6 & 12 & \href{https://www.ncbi.nlm.nih.gov/geo/query/acc.cgi?acc=GSE22809}{GSE22809} & \href{https://doi.org/10.1371/journal.pgen.1003128}{doi} \\
\addlinespace[3pt]
306 & \textbf{Rossmann 2011}\newline \sourcetext{A common telomeric gene silencing assay is affected by nucleotide metabolism. \emph{Mol Cell}} 2011. & 1 & -- & 1 & 6 & 12 & \href{https://www.ncbi.nlm.nih.gov/geo/query/acc.cgi?acc=GSE27222}{GSE27222} & \href{https://doi.org/10.1016/j.molcel.2011.03.007}{doi} \\
\addlinespace[3pt]
165 & \textbf{Sapra 2004}\newline \sourcetext{Genome-wide analysis of pre-mRNA splicing: intron features govern the requirement for the second-step factor, Prp17 in Saccharomyces cerevisiae and Schizosaccharomyces pombe. \emph{J Biol Chem}} 2004. & 1 & -- & 1 & 6 & 24 & \textbf{none} & \href{https://doi.org/10.1074/jbc.M408815200}{doi} \\
\addlinespace[3pt]
205 & \textbf{Shivaswamy 2008}\newline \sourcetext{Stress-dependent dynamics of global chromatin remodeling in yeast: dual role for SWI/SNF in the heat shock stress response. \emph{Mol Cell Biol}} 2008. & 1 & -- & 1 & 6 & 19 & \href{https://www.ncbi.nlm.nih.gov/geo/query/acc.cgi?acc=GSE7665}{GSE7665} & \href{https://doi.org/10.1128/MCB.01659-07}{doi} \\
\addlinespace[3pt]
369 & \textbf{Dong 2007}\newline \sourcetext{YRA1 autoregulation requires nuclear export and cytoplasmic Edc3p-mediated degradation of its pre-mRNA. \emph{Mol Cell}} 2007. & 1 & -- & 1 & 5 & 10 & \href{https://www.ncbi.nlm.nih.gov/geo/query/acc.cgi?acc=GSE6647}{GSE6647} & \href{https://doi.org/10.1016/j.molcel.2007.01.012}{doi} \\
\addlinespace[3pt]
164 & \textbf{Aristizabal 2013}\newline \sourcetext{High-throughput genetic and gene expression analysis of the RNAPII-CTD reveals unexpected connections to SRB10/CDK8. \emph{PLoS Genet}} 2013. & 1 & -- & 1 & 4 & 24 & \href{https://www.ncbi.nlm.nih.gov/geo/query/acc.cgi?acc=GSE43120}{GSE43120} & \href{https://doi.org/10.1371/journal.pgen.1003758}{doi} \\
\addlinespace[3pt]
300 & \textbf{Chen 2006}\newline \sourcetext{Coordinate regulation of multiple and distinct biosynthetic pathways by TOR and PKA kinases in S. cerevisiae. \emph{Curr Genet}} 2006. & 1 & -- & 1 & 4 & 12 & \href{https://www.ncbi.nlm.nih.gov/geo/query/acc.cgi?acc=GSE3160}{GSE3160} & \href{https://doi.org/10.1007/s00294-005-0055-9}{doi} \\
\addlinespace[3pt]
302 & \textbf{Drozdova 2014}\newline \sourcetext{Transcriptional response to the [ISP(+) ] prion of Saccharomyces cerevisiae differs from that induced by the deletion of its structural gene, SFP1. \emph{FEMS Yeast Res}} 2014. & 1 & -- & 1 & 4 & 12 & \href{https://www.ncbi.nlm.nih.gov/geo/query/acc.cgi?acc=GSE52189}{GSE52189} & \href{https://doi.org/10.1111/1567-1364.12211}{doi} \\
\addlinespace[3pt]
409 & \textbf{Klopf 2017}\newline \sourcetext{INO80 represses osmostress induced gene expression by resetting promoter proximal nucleosomes. \emph{Nucleic Acids Res}} 2017. & 1 & -- & 1 & 4 & 8 & \href{https://www.ncbi.nlm.nih.gov/geo/query/acc.cgi?acc=GSE78766}{GSE78766} & \href{https://doi.org/10.1093/nar/gkw1292}{doi} \\
\addlinespace[3pt]
36 & \textbf{Lee 2005}\newline \sourcetext{Multiple RNA surveillance pathways limit aberrant expression of iron uptake mRNAs and prevent iron toxicity in S. cerevisiae. \emph{Mol Cell}} 2005. & 1 & 1 & 1 & 4 & 9 & \textbf{none} & \href{https://doi.org/10.1016/j.molcel.2005.05.021}{doi} \\
\addlinespace[3pt]
410 & \textbf{Liko 2007}\newline \sourcetext{Stb3 binds to ribosomal RNA processing element motifs that control transcriptional responses to growth in Saccharomyces cerevisiae. \emph{J Biol Chem}} 2007. & 1 & -- & 1 & 4 & 8 & \href{https://www.ncbi.nlm.nih.gov/geo/query/acc.cgi?acc=GSE8379}{GSE8379} & \href{https://doi.org/10.1074/jbc.M704762200}{doi} \\
\addlinespace[3pt]
35 & \textbf{Alkim 2013}\newline \sourcetext{Mechanisms other than activation of the iron regulon account for the hyper-resistance to cobalt of a Saccharomyces cerevisiae strain obtained by evolutionary engineering. \emph{Metallomics}} 2013. & 1 & 1 & 1 & 3 & 9 & \href{https://www.ncbi.nlm.nih.gov/geo/query/acc.cgi?acc=GSE39185}{GSE39185} & \href{https://doi.org/10.1039/c3mt00107e}{doi} \\
\addlinespace[3pt]
408 & \textbf{Angus-Hill 2001}\newline \sourcetext{A Rsc3/Rsc30 zinc cluster dimer reveals novel roles for the chromatin remodeler RSC in gene expression and cell cycle control. \emph{Mol Cell}} 2001. & 1 & -- & 1 & 3 & 8 & \textbf{none} & \href{https://doi.org/10.1016/s1097-2765(01)00219-2}{doi} \\
\addlinespace[3pt]
275 & \textbf{Arana 2012}\newline \sourcetext{Transcriptional responses to loss of RNase H2 in Saccharomyces cerevisiae. \emph{DNA Repair (Amst)}} 2012. & 1 & -- & 1 & 3 & 14 & \href{https://www.ncbi.nlm.nih.gov/geo/query/acc.cgi?acc=GSE34668}{GSE34668} & \href{https://doi.org/10.1016/j.dnarep.2012.09.006}{doi} \\
\addlinespace[3pt]
451 & \textbf{Chia 2018}\newline \sourcetext{Knockout of the Hmt1p Arginine Methyltransferase in Saccharomyces cerevisiae Leads to the Dysregulation of Phosphate-associated Genes and Processes. \emph{Mol Cell Proteomics}} 2018. & 1 & -- & 1 & 3 & 7 & \href{https://www.ncbi.nlm.nih.gov/geo/query/acc.cgi?acc=GSE99869}{GSE99869} & \href{https://doi.org/10.1074/mcp.RA117.000214}{doi} \\
\addlinespace[3pt]
40 & \textbf{Eriksson 2005}\newline \sourcetext{Global regulation by the yeast Spt10 protein is mediated through chromatin structure and the histone upstream activating sequence elements. \emph{Mol Cell Biol}} 2005. & 1 & 1 & 1 & 3 & 6 & \textbf{none} & \href{https://doi.org/10.1128/MCB.25.20.9127-9137.2005}{doi} \\
\addlinespace[3pt]
42 & \textbf{Garca-Rodrguez 2015}\newline \sourcetext{Manganese redistribution by calcium-stimulated vesicle trafficking bypasses the need for P-type ATPase function. \emph{J Biol Chem}} 2015. & 1 & 1 & 1 & 3 & 6 & \href{https://www.ncbi.nlm.nih.gov/geo/query/acc.cgi?acc=GSE62161}{GSE62161} & \href{https://doi.org/10.1074/jbc.M114.616334}{doi} \\
\addlinespace[3pt]
43 & \textbf{Garcia-Rodriguez 2012}\newline \sourcetext{Impaired manganese metabolism causes mitotic misregulation. \emph{J Biol Chem}} 2012. & 1 & 1 & 1 & 3 & 6 & \href{https://www.ncbi.nlm.nih.gov/geo/query/acc.cgi?acc=GSE29420}{GSE29420} & \href{https://doi.org/10.1074/jbc.M112.358309}{doi} \\
\addlinespace[3pt]
44 & \textbf{Gavalda 2016}\newline \sourcetext{Excess of Yra1 RNA-Binding Factor Causes Transcription-Dependent Genome Instability, Replication Impairment and Telomere Shortening. \emph{PLoS Genet}} 2016. & 1 & 1 & 1 & 3 & 6 & \href{https://www.ncbi.nlm.nih.gov/geo/query/acc.cgi?acc=GSE68487}{GSE68487} & \href{https://doi.org/10.1371/journal.pgen.1005966}{doi} \\
\addlinespace[3pt]
465 & \textbf{Hickman 2007}\newline \sourcetext{Heme levels switch the function of Hap1 of Saccharomyces cerevisiae between transcriptional activator and transcriptional repressor. \emph{Mol Cell Biol}} 2007. & 1 & -- & 1 & 3 & 6 & \href{https://www.ncbi.nlm.nih.gov/geo/query/acc.cgi?acc=GSE8613}{GSE8613} & \href{https://doi.org/10.1128/MCB.00887-07}{doi} \\
\addlinespace[3pt]
34 & \textbf{Hopfler 2019}\newline \sourcetext{Slx5/Slx8-dependent ubiquitin hotspots on chromatin contribute to stress tolerance. \emph{EMBO J}} 2019. & 1 & 1 & 1 & 3 & 12 & \href{https://www.ncbi.nlm.nih.gov/geo/query/acc.cgi?acc=GSE118817}{GSE118817} & \href{https://doi.org/10.15252/embj.2018100368}{doi} \\
\addlinespace[3pt]
261 & \textbf{Lanza 2012}\newline \sourcetext{Linking yeast Gcn5p catalytic function and gene regulation using a quantitative, graded dominant mutant approach. \emph{PLoS One}} 2012. & 1 & -- & 1 & 3 & 15 & \href{https://www.ncbi.nlm.nih.gov/geo/query/acc.cgi?acc=GSE26923}{GSE26923} & \href{https://doi.org/10.1371/journal.pone.0036193}{doi} \\
\addlinespace[3pt]
37 & \textbf{Mendiratta 2006}\newline \sourcetext{The DNA-binding domain of the yeast Spt10p activator includes a zinc finger that is homologous to foamy virus integrase. \emph{J Biol Chem}} 2006. & 1 & 1 & 1 & 3 & 9 & \href{https://www.ncbi.nlm.nih.gov/geo/query/acc.cgi?acc=GSE3354}{GSE3354} & \href{https://doi.org/10.1074/jbc.M511416200}{doi} \\
\addlinespace[3pt]
466 & \textbf{Rodriguez-Navarro 2004}\newline \sourcetext{Sus1, a functional component of the SAGA histone acetylase complex and the nuclear pore-associated mRNA export machinery. \emph{Cell}} 2004. & 1 & -- & 1 & 3 & 6 & \textbf{none} & \href{https://doi.org/10.1016/s0092-8674(03)01025-0}{doi} \\
\addlinespace[3pt]
467 & \textbf{Santos-Pereira 2013}\newline \sourcetext{The Npl3 hnRNP prevents R-loop-mediated transcription-replication conflicts and genome instability. \emph{Genes Dev}} 2013. & 1 & -- & 1 & 3 & 6 & \href{https://www.ncbi.nlm.nih.gov/geo/query/acc.cgi?acc=GSE50186}{GSE50186} & \href{https://doi.org/10.1101/gad.229880.113}{doi} \\
\addlinespace[3pt]
588 & \textbf{Costelloe 2012}\newline \sourcetext{The yeast Fun30 and human SMARCAD1 chromatin remodellers promote DNA end resection. \emph{Nature}} 2012. & 1 & -- & 1 & 2 & 2 & \href{https://www.ncbi.nlm.nih.gov/geo/query/acc.cgi?acc=GSE38735}{GSE38735} & \href{https://doi.org/10.1038/nature11353}{doi} \\
\addlinespace[3pt]
26 & \textbf{Garcia-Martinez 2016}\newline \sourcetext{The cellular growth rate controls overall mRNA turnover, and modulates either transcription or degradation rates of particular gene regulons. \emph{Nucleic Acids Res}} 2016. & 1 & 1 & 1 & 2 & 32 & \href{https://www.ncbi.nlm.nih.gov/geo/query/acc.cgi?acc=GSE63729}{GSE63729}\newline \href{https://www.ncbi.nlm.nih.gov/geo/query/acc.cgi?acc=GSE57467}{GSE57467}\newline \href{https://www.ncbi.nlm.nih.gov/geo/query/acc.cgi?acc=GSE65225}{GSE65225} & \href{https://doi.org/10.1093/nar/gkv1512}{doi} \\
\addlinespace[3pt]
48 & \textbf{Kang 2014}\newline \sourcetext{A novel protein, Pho92, has a conserved YTH domain and regulates phosphate metabolism by decreasing the mRNA stability of PHO4 in Saccharomyces cerevisiae. \emph{Biochem J}} 2014. & 1 & 1 & 1 & 2 & 2 & \href{https://www.ncbi.nlm.nih.gov/geo/query/acc.cgi?acc=GSE50086}{GSE50086} & \href{https://doi.org/10.1042/BJ20130862}{doi} \\
\addlinespace[3pt]
576 & \textbf{Levy 2008}\newline \sourcetext{Yeast linker histone Hho1p is required for efficient RNA polymerase I processivity and transcriptional silencing at the ribosomal DNA. \emph{Proc Natl Acad Sci U S A}} 2008. & 1 & -- & 1 & 2 & 3 & \href{https://www.ncbi.nlm.nih.gov/geo/query/acc.cgi?acc=GSE9095}{GSE9095} & \href{https://doi.org/10.1073/pnas.0709403105}{doi} \\
\addlinespace[3pt]
526 & \textbf{Lu 2005}\newline \sourcetext{The YJR127C/ZMS1 gene product is involved in glycerol-based respiratory growth of the yeast Saccharomyces cerevisiae. \emph{Curr Genet}} 2005. & 1 & -- & 1 & 2 & 5 & \href{https://www.ncbi.nlm.nih.gov/geo/query/acc.cgi?acc=GSE29530}{GSE29530} & \href{https://doi.org/10.1007/s00294-005-0023-4}{doi} \\
\addlinespace[3pt]
370 & \textbf{Lu 2022}\newline \sourcetext{A balancing act: interactions within NuA4/TIP60 regulate picNuA4 function in Saccharomyces cerevisiae and humans. \emph{Genetics}} 2022. & 1 & -- & 1 & 2 & 10 & \href{https://www.ncbi.nlm.nih.gov/geo/query/acc.cgi?acc=GSE137261}{GSE137261} & \href{https://doi.org/10.1093/genetics/iyac136}{doi} \\
\addlinespace[3pt]
589 & \textbf{Radman-Livaja 2012}\newline \sourcetext{A key role for Chd1 in histone H3 dynamics at the 3' ends of long genes in yeast. \emph{PLoS Genet}} 2012. & 1 & -- & 1 & 2 & 2 & \href{https://www.ncbi.nlm.nih.gov/geo/query/acc.cgi?acc=GSE38496}{GSE38496} & \href{https://doi.org/10.1371/journal.pgen.1002811}{doi} \\
\addlinespace[3pt]
305 & \textbf{Rodriguez-Colman 2010}\newline \sourcetext{The forkhead transcription factor Hcm1 promotes mitochondrial biogenesis and stress resistance in yeast. \emph{J Biol Chem}} 2010. & 1 & -- & 1 & 2 & 12 & \href{https://www.ncbi.nlm.nih.gov/geo/query/acc.cgi?acc=GSE20420}{GSE20420} & \href{https://doi.org/10.1074/jbc.M110.174763}{doi} \\
\addlinespace[3pt]
45 & \textbf{Rudra 2005}\newline \sourcetext{Central role of Ifh1p-Fhl1p interaction in the synthesis of yeast ribosomal proteins. \emph{EMBO J}} 2005. & 1 & 1 & 1 & 2 & 6 & \textbf{none} & \href{https://doi.org/10.1038/sj.emboj.7600553}{doi} \\
\addlinespace[3pt]
209 & \textbf{Sabet 2004}\newline \sourcetext{Genome-wide analysis of the relationship between transcriptional regulation by Rpd3p and the histone H3 and H4 amino termini in budding yeast. \emph{Mol Cell Biol}} 2004. & 1 & -- & 1 & 2 & 18 & \textbf{none} & \href{https://doi.org/10.1128/MCB.24.20.8823-8833.2004}{doi} \\
\addlinespace[3pt]
276 & \textbf{Saint 2014}\newline \sourcetext{The TAF9 C-terminal conserved region domain is required for SAGA and TFIID promoter occupancy to promote transcriptional activation. \emph{Mol Cell Biol}} 2014. & 1 & -- & 1 & 2 & 14 & \href{https://www.ncbi.nlm.nih.gov/geo/query/acc.cgi?acc=GSE44544}{GSE44544} & \href{https://doi.org/10.1128/MCB.01060-13}{doi} \\
\addlinespace[3pt]
47 & \textbf{Sariki 2016}\newline \sourcetext{Sen1, the homolog of human Senataxin, is critical for cell survival through regulation of redox homeostasis, mitochondrial function, and the TOR pathway in Saccharomyces cerevisiae. \emph{FEBS J}} 2016. & 1 & 1 & 1 & 2 & 4 & \href{https://www.ncbi.nlm.nih.gov/geo/query/acc.cgi?acc=GSE75447}{GSE75447} & \href{https://doi.org/10.1111/febs.13917}{doi} \\
\addlinespace[3pt]
599 & \textbf{Davis 2006}\newline \sourcetext{Accumulation of unstable promoter-associated transcripts upon loss of the nuclear exosome subunit Rrp6p in Saccharomyces cerevisiae. \emph{Proc Natl Acad Sci U S A}} 2006. & 1 & -- & 1 & 1 & 1 & \href{https://www.ncbi.nlm.nih.gov/geo/query/acc.cgi?acc=GSE3813}{GSE3813} & \href{https://doi.org/10.1073/pnas.0507783103}{doi} \\
\addlinespace[3pt]
242 & \textbf{Hughes Hallett 2014}\newline \sourcetext{State transitions in the TORC1 signaling pathway and information processing in Saccharomyces cerevisiae. \emph{Genetics}} 2014. & 1 & -- & 1 & 1 & 16 & \href{https://www.ncbi.nlm.nih.gov/geo/query/acc.cgi?acc=GSE58992}{GSE58992} & \href{https://doi.org/10.1534/genetics.114.168369}{doi} \\
\addlinespace[3pt]
535 & \textbf{Sen 2011}\newline \sourcetext{A new, highly conserved domain in Swi2/Snf2 is required for SWI/SNF remodeling. \emph{Nucleic Acids Res}} 2011. & 1 & -- & 1 & 1 & 4 & \href{https://www.ncbi.nlm.nih.gov/geo/query/acc.cgi?acc=GSE30675}{GSE30675} & \href{https://doi.org/10.1093/nar/gkr622}{doi} \\
\end{longtable}
\endgroup
\end{landscape}


\subsection{Conditions labeled by gene name alone\secstatus{todo}}\label{sec:gene-named}

\spellNNamedStudies{} studies label conditions with nothing but gene names, for
\spellNNamedConditions{} conditions and \spellNNamedGenes{} distinct genes
(Table~\ref{tab:gene-named}). These name a perturbed gene without saying how it
was perturbed, so they are kept out of the deletion count. Two studies hold
almost all of them. The headers of Hughes 2000 mix lowercase names with entries
marked \texttt{(tet promoter)}, and those of Mnaimneh 2004 pair a standard name
with its systematic name. Hypothesis (untested against the papers): the
lowercase Hughes 2000 names are deletion mutants and the Mnaimneh 2004 entries
are tetracycline-repressible promoter alleles of essential genes.

%% GENERATED FILE -- do not hand-edit.
%% SOURCE: experiments/015-spell/scripts/spell_publication_tables.py
\begin{table}[H]\centering
\footnotesize
\caption[]{Studies with conditions labeled by gene name alone. \emph{\#} is the
row number in Table~\ref{tab:all}. \emph{genes} counts the distinct genes named,
\emph{new} the gene-named genotypes in neither Kemmeren 2014 nor Sameith 2015,
\emph{named cond.} the conditions labeled that way and \emph{cond.} the study's
total. The header does not say how the gene was perturbed, so these are not
counted as deletions.}\label{tab:gene-named}
\begin{tabular}{rlrrrrll}
\toprule
\# & study & genes & new & named cond. & cond. & GEO & link \\
\midrule
2 & \textbf{Hughes 2000} & 276 & 157 & 284 & 300 & \textbf{none} & \href{https://doi.org/10.1016/s0092-8674(00)00015-5}{doi} \\
1 & \textbf{Mnaimneh 2004} & 215 & 214 & 215 & 215 & \textbf{none} & \href{https://doi.org/10.1016/j.cell.2004.06.013}{doi} \\
212 & \textbf{Bernstein 2000} & 3 & -- & 6 & 18 & \textbf{none} & \href{https://doi.org/10.1073/pnas.250477697}{doi} \\
176 & \textbf{Meneghini 2003} & 3 & -- & 12 & 24 & \textbf{none} & \href{https://doi.org/10.1016/s0092-8674(03)00123-5}{doi} \\
17 & \textbf{Hong 2012} & 2 & 2 & 2 & 14 & \href{https://www.ncbi.nlm.nih.gov/geo/query/acc.cgi?acc=GSE36118}{GSE36118} & \href{https://doi.org/10.1128/AEM.01444-12}{doi} \\
14 & \textbf{Zheng 2010} & 2 & 2 & 2 & 48 & \href{https://www.ncbi.nlm.nih.gov/geo/query/acc.cgi?acc=GSE19634}{GSE19634} & \href{https://doi.org/10.1038/nature08934}{doi} \\
384 & \textbf{Travers 2000} & 1 & -- & 1 & 10 & \textbf{none} & \href{https://doi.org/10.1016/s0092-8674(00)80835-1}{doi} \\
\bottomrule
\end{tabular}
\end{table}

